% from ega3-1-fr.tex
\oldpage[III]{122}
et comme le faisceau $R^nf_*(\mathcal F)$ est associé au préfaisceau
$\mathfrak Y_g\mapsto H^n(f^{-1}(\mathfrak Y_g),\mathcal F)$ sur les
$\mathfrak Y_g$ (0\textsubscript{I}, 3.2.1), on a bien montré que
l'homomorphisme (\hyperref[III.3.4.2.2-fr]{3.4.2.2}) est bijectif. Prouvons
ensuite que $R^nf_*(\mathcal F)$ est un $\mathcal O_{\mathfrak Y}$-Module
cohérent, et de façon plus précise que l'on a
\[
R^nf_*(\mathcal F)=(H^n(\mathfrak X,\mathcal F))^\Delta.
\tag{3.4.5.1}\label{III.3.4.5.1-fr}
\]
Avec les notations précédentes, on a, puisque $\mathcal F_k$ est un
$\mathcal O_{X_k}$-Module cohérent (\hyperref[III.1.4.13-fr]{1.4.13})
\[
\Gamma(\mathfrak Y_g,R^nf_*(\mathcal F_k))
=(\Gamma(\mathfrak Y,R^nf_*(\mathcal F_k)))_g
=(H^n(\mathfrak X,\mathcal F_k))_g.
\]
Or, les $H^n(\mathfrak X,\mathcal F_k)$ forment un système projectif
vérifiant (ML) et leur limite projective $H^n(\mathfrak X,\mathcal F)$ est
un $A$-module de type fini. On en conclut (0, 13.7.8) que l'on a
\[
\varprojlim_k(H^n(\mathfrak X,\mathcal F_k))_g
=H^n(\mathfrak X,\mathcal F)\otimes_AA_{\{g\}}
=\Gamma(\mathfrak Y_g,(H^n(\mathfrak X,\mathcal F))^\Delta)
\]
compte tenu de (I, 10.10.8) appliqué à $A$ et $A_{\{g\}}$~; ceci démontre
(\hyperref[III.3.4.5.1-fr]{3.4.5.1}) puisque
\[
\Gamma(\mathfrak Y_g,R^nf_*(\mathcal F))
=\varprojlim_k\Gamma(\mathfrak Y_g,R^nf_*(\mathcal F_k)).
\]
Comme (\hyperref[III.3.4.2.2-fr]{3.4.2.2}) est alors un isomorphisme de
$\mathcal O_{\mathfrak Y}$-Modules cohérents, c'est nécessairement un
isomorphisme \emph{topologique} (I, 10.11.6). Enfin, il résulte des relations
$R^nf_*(\mathcal F_k)=(H^n(\mathfrak X,\mathcal F_k))^\Delta$ que le système
projectif $(R^nf_*(\mathcal F_k))_{k\geq0}$ vérifie (ML) (I, 10.10.2).

Une fois (\hyperref[III.3.4.2-fr]{3.4.2}) et
(\hyperref[III.3.4.3-fr]{3.4.3}) démontrés dans le cas où le préschéma formel
$\mathfrak Y$ est affine noethérien, il est immédiat de passer de là au cas
général pour (\hyperref[III.3.4.2-fr]{3.4.2}) et la première assertion de
(\hyperref[III.3.4.3-fr]{3.4.3}), qui sont locales sur $\mathfrak Y$. Quant à
la seconde assertion de (\hyperref[III.3.4.3-fr]{3.4.3}), il suffit,
$\mathfrak Y$ étant noethérien, de le recouvrir par un nombre fini d'ouverts
affines noethériens $U_i$ et de remarquer que les restrictions du système
projectif $(R^qf_*(\mathcal F_k))$ à chacun des $U_i$ vérifient (ML).
\end{env}

Nous avons en outre démontré en cours de route~:

\begin{corollary}[3.4.6]
\label{III.3.4.6-fr}
Sous les hypothèses de (\hyperref[III.3.4.4-fr]{3.4.4}), l'homomorphisme
canonique
\[
H^q(\mathfrak X,\mathcal F)
\longrightarrow\Gamma(\mathfrak Y,R^qf_*(\mathcal F))
\tag{3.4.6.1}\label{III.3.4.6.1-fr}
\]
est bijectif.
\end{corollary}

\section{Le théorème fondamental des morphismes propres. Applications}
\label{section:III.4-fr}

\subsection{Le théorème fondamental}
\label{subsection:III.4.1-fr}

\begin{env}[4.1.1]
\label{III.4.1.1-fr}
Soient $X$, $Y$ deux préschémas usuels noethériens, $f:X\to Y$ un
morphisme propre, $Y'$ une partie fermée de $Y$, $X'$ son image réciproque
$f^{-1}(Y')$. Nous désignerons par $\widehat X$ et $\widehat Y$ les
préschémas formels $X_{/X'}$ et $Y_{/Y'}$, complétés de $X$ et $Y$ le long
de $X'$ et $Y'$ respectivement (I, 10.8.5), par $\widehat f$ le prolongement
de $f$ à ces complétés (I, 10.9.1) qui est un morphisme
$\widehat X\to\widehat Y$ de préschémas formels. Pour tout
$\mathcal O_X$-Module cohérent $\mathcal F$, nous

